\uuid{V7u7}
\exo7id{4790}
\titre{exo7 4790}
\auteur{quercia}
\organisation{exo7}
\datecreate{2010-03-16}
\isIndication{false}
\isCorrection{false}
\chapitre{Topologie}
\sousChapitre{Topologie des espaces vectoriels normés}
\cpge{mpsi-mp2i.limites-continuite.b, pcsi.limites-continuite.b, ptsi.limites-continuite.b, mpsi-mp2i.derivabilite.a, pcsi.derivabilite.a, ptsi.derivabilite.a, mpsi-mp2i.fonctions-deux-variables.b, pcsi.fonctions-deux-variables.b, ptsi.fonctions-deux-variables.b, bcpst1.limites-continuite.c, bcpst1.derivation.a, mp-mpi.topologie-espaces-normes.e, pc.espaces-vectoriels-normes.e, psi.espaces-vectoriels-normes.e, pt.fonctions-deux-trois-variables.b}
\module{Analyse}
\niveau{L2}
\difficulte{}

\contenu{
\question{
Soit $f:\R \to \R$ une fonction de classe $\mathcal{C}^1$ et
$g : {\R^2} \to \R$ $\begin{cases}{(x,y)} &{\frac {f(x)-f(y)}{x-y} \text{ si } x \ne y}\cr
 {(x,x)} &{f'(x)}\cr\end{cases}$

Montrer que $g$ est continue.
(Attention : pour une fonction d{\'e}finie par cas, se placer au voisinage d'un
point $(x_0,y_0)$ et d{\'e}terminer si un seul ou plusieurs cas sont {\`a} consid{\'e}rer
dans ce voisinage)
}
}
